\section{The Act Under Study}

On 8 December 2010, the Solemnity of the Immaculate Conception, David
Laurin Ricken, twelfth Bishop of Green Bay, signed a four-page decree
in the chapel at Champion, Wisconsin. Its title states its object
precisely: \work{Decree on the Authenticity of the Apparitions of 1859
at the Shrine of Our Lady of Good Help}. Everything this study says
about what the Catholic Church has and has not settled at Champion
hangs on reading the last page of that decree exactly.

\begin{operativetext}{Decree of David L. Ricken, Bishop of Green Bay, 8
December 2010, operative clause, read in the four-page signed PDF
published by the National Shrine}
It remains to me now, the Twelfth Bishop of the Diocese of Green Bay
and the lowliest of the servants of Mary, to declare with moral
certainty and in accord with the norms of the Church:\par
\textbf{that the events, apparitions and locutions given to Adele Brise
in October of 1859 do exhibit the substance of supernatural character,
and I do hereby approve these apparitions as worthy of belief (although
not obligatory) by the Christian faithful.}\par
I encourage the faithful to frequent this holy place as a place of
solace and answered prayer.\par
Given at the Shrine of Our Lady of Good Help, Champion, Wisconsin, the
eighth day of December in the year of Our Lord two thousand and ten,
the Solemnity of the Immaculate Conception of the Blessed Virgin
Mary.\endnote{David L. Ricken, Bishop of Green Bay, \work{Decree on the
Authenticity of the Apparitions of 1859 at the Shrine of Our Lady of
Good Help, Diocese of Green Bay}, 8 December 2010, p.~4. Read
2026-07-25 in the four-page PDF published by the National Shrine of Our
Lady of Champion at
\url{https://championshrine.org/wp-content/uploads/2022/06/Shrine-of-Our-Lady-of-Good-Help-Declaration.pdf}
(SHA-256 of the file as retrieved:
\texttt{b10ad96f0821fef8072164d6f83d3f33d9e008e1aa510d1c5706bbe53bbd664e}).
Pages~1--3 carry an extractable text layer; page~4, which bears the
operative clause and the signatures, is an image, and its wording here
was read directly from a 150~dpi rendering of that page. The bold
emphasis on the operative sentence is the decree's own. The page bears
the signature of Bishop Ricken, the countersignature of the Chancellor,
John F. Doerfler, and the diocesan seal with the attestation
\term{Concordat cum Originali}. The diocesan original, and any protocol
number or register entry, were not consulted (Appendix~A). All web loci
cited in these notes were read on 2026-07-25 unless another date is
stated.}
\end{operativetext}

Four features of that sentence carry the whole weight of the case, and
each of them is easy to lose in summary.

First, the \emph{object}. What is approved is ``the events, apparitions
and locutions given to Adele Brise in October of 1859.'' Not a shrine,
not a devotion, not a book, not a fire, not a healing, not a
chronology, not a text. The decree fixes a month and a year and a
person, and it fixes nothing else about the events themselves --- not a
day, not a number of appearances, not a sequence, not a wording.

Second, the \emph{predicate}. The events ``do exhibit the substance of
supernatural character.'' That is not the older formula \term{constat
de supernaturalitate}. ``The substance of'' is a real qualifier, and
the decree's own reasoning, quoted in \S2.2 below, shows why the bishop
needed it: his experts were working from a documentary record he
himself describes as ``not abundant.'' The claim is that what is
substantially there is of supernatural character, not that every
transmitted detail is.

Third, the \emph{effect}. The apparitions are approved ``as worthy of
belief (although not obligatory) by the Christian faithful.'' The
parenthesis is in the decree. It is the whole doctrine of private
revelation compressed into three words: public Revelation is complete
in Christ, and no approved apparition can add to it or require the
divine and Catholic faith owed to it.\endnote{Second Vatican Council,
\work{Dei Verbum} (18 November 1965), n.~4, quoted in \work{Catechism
of the Catholic Church}, n.~66: ``no new public revelation is to be
expected before the glorious manifestation of our Lord Jesus Christ.''
For the standing account of what ecclesiastical approval of a private
revelation means, see Benedict XVI, \work{Verbum Domini} (30 September
2010), n.~14, quoted in full at note~24 of the Dicastery for the
Doctrine of the Faith, \work{Norms for Proceeding in the Discernment of
Alleged Supernatural Phenomena} (17 May 2024): ``Ecclesiastical
approval of a private revelation essentially means that its message
contains nothing contrary to faith and morals; it is licit to make it
public and the faithful are authorized to give to it their adherence in
a prudent manner. \dots{} It is a help which is proffered, but its use
is not obligatory.'' English text of the Norms read at
\url{https://www.vatican.va/roman_curia/congregations/cfaith/documents/rc_ddf_doc_20240517_norme-fenomeni-soprannaturali_en.html}.}

Fourth, the \emph{author}. This is a diocesan act. It was made by the
local ordinary under the procedural regime then in force, the
Congregation for the Doctrine of the Faith's \work{Norms regarding the
manner of proceeding in the discernment of presumed apparitions or
revelations} of 25~February 1978, whose Part~III begins: ``Above all,
the duty of vigilance and intervention falls to the Ordinary of the
place.''\endnote{Sacred Congregation for the Doctrine of the Faith,
\work{Norms regarding the manner of proceeding in the discernment of
presumed apparitions or revelations} (25 February 1978), III.1. English
text read at
\url{https://www.vatican.va/roman_curia/congregations/cfaith/documents/rc_con_cfaith_doc_19780225_norme-apparizioni_en.html}.
The Norms were approved by Paul~VI on 24 February 1978 and dated 25
February 1978; they were signed by Cardinal \v{S}eper as Prefect and
J\'er\^ome Hamer, O.P., as Secretary. This study makes no claim about
when or in what form the text was first made publicly available; see
Appendix~A.} The decree names no Roman intervention, and \S7 explains
why its silence on that point settles nothing either way.

\studythesis{This is an evidence-first apparition study. Its governing
claim is that Champion is the clearest American instance of a
positively judged private revelation whose event stratum has almost no
contemporary documentation, and that this is a statement about the
record rather than an argument against the judgment. The two earliest
published narratives that this study could reach --- Eliza Allen
Starr's, copyright-entered in 1871, and Peter Pernin's, printed at
Montreal in 1874 --- were both written more than a decade after the
events, both by devotional authors raising money, and they disagree
with each other and with the current custodial narrative about the
season, the distance walked, the order and setting of the encounters,
the appearance of the figure, and above all the extent of the reported
words. The self-identification as Queen of Heaven, the command of
general confession, the warning of punishment, the beatitude on the
companions, and the question about the vineyard are absent from both
of them. The preservation of the chapel enclosure in the fires of 8
October 1871 is early and well attested by Pernin; the downpour, the
refugee crowd, the livestock, and the well are not, and a widely
reproduced paragraph about that night circulates under Pernin's name
that he did not write. Bishop Ricken's decree of 2010 judged, with
moral certainty, that the October 1859 events exhibit the substance of
supernatural character and approved them as worthy of belief though not
obligatory. That judgment deserves exact reception. It does not convert
any later sentence into a transcript, does not declare the fire or any
cure a miracle --- the decree says in terms that none has been so
declared --- and does not add anything to the Revelation complete in
Christ.}

The study runs in the order the evidence forces. \S2 states what the
record does not contain, because at Champion the absence is the
governing fact. \S3 and \S4 read the two earliest published witnesses
at length and in their own settings. \S5 collates them against each
other and against the narrative the shrine publishes today, and sets
out the stratigraphy of the reported message. \S6 separates the fire of
8 October 1871 into what contemporary record supports and what later
devotional narration supplies, and documents a misattributed quotation
that has passed into general circulation. \S7 reads the 2010 inquiry
and decree against the norms then in force and states exactly what the
act obliges. \S8 follows the later institutional acts --- national
shrine, change of title, annual solemnity, cause of canonization ---
and asks what the 2024 norms did to a legacy judgment of this kind.
\S9 states this study's own judgment under a visible label, with the
strongest counterargument and the finding that would overturn it.
Work-wide bounds, taxonomy, and qualifications are in the terminal
appendix; witnesses and rights are in the References.

\begin{statusdossier}{Controlling status at the as-of date of this
study (2026-07-25)}
\dosfield{Event}{Reported apparitions and locutions to Adele Brise,
Robinsonville (now Champion), Brown County, Wisconsin, October 1859.}
\dosfield{Competent authority}{The Bishop of Green Bay. The see was not
erected until 1868; the jurisdiction competent in 1859 is not named in
this study, which did not verify it.}
\dosfield{Controlling act}{Decree of Bishop David L. Ricken, 8 December
2010, approving the October 1859 events, apparitions, and locutions as
exhibiting the substance of supernatural character and as worthy of
belief although not obligatory.}
\dosfield{Procedural regime of that act}{The CDF Norms of 25 February
1978, replaced in their entirety on 19 May 2024 by the Norms of the
Dicastery for the Doctrine of the Faith of 17 May 2024 (art.~27 of the
latter).}
\dosfield{What the act does not establish}{Any particular day, count,
sequence, or wording of the 1859 events; the authenticity of any
narrative detail beyond ``the substance''; a miracle, whether of the
1871 preservation or of any reported healing --- the decree states that
none of the reported favours has been officially declared a miracle;
any obligation of belief.}
\dosfield{Later acts located}{Designation as a national shrine
(reported by the shrine as 19 March 2016, by the United States
Conference of Catholic Bishops; the conference's own act was not
reached); a change of title from Our Lady of Good Help to Our Lady of
Champion and an annual solemnity on 9 October, first celebrated in
2023, which the shrine attributes to word received from the Holy See's
Dicastery for Divine Worship in 2023 (the underlying instruments were
not reached); the opening of a cause of canonization for Adele Brise,
who is now styled Servant of God.}
\dosfield{Not located}{Any act of the Dicastery for the Doctrine of the
Faith, before or after 2010, reclassifying, confirming, restricting, or
superseding the 2010 decree.}
\end{statusdossier}
